%==============================================================================
% Section 7: Physical Interpretation and Outlook
%==============================================================================
\section{Physical Interpretation and Outlook}
\label{sec:discussion}

\subsection{Physical Hierarchy of the Three Topologies}
\label{sec:discussion:hierarchy}

Through the analysis of this paper, the three topologies
$\Sthree\times\Sone$, $\Tthree\times\Sone$, and $\Nilthree\times\Sone$
emerge as having sharply distinct physical roles with respect to the EC--Weyl
coupling.

\paragraph{$\Tthree\times\Sone$.}
As a flat manifold ($C^a_{bc}=0$), all spin-2/spin-1 mode masses are zero, and
the EC--Weyl coupling generates no new stable point in any of the AX/VT/MX
backgrounds (triple triviality, \S\ref{sec:t3:triple}).  $\Tthree$ functions as
the ``trivial reference point'' of the theory, with the asymptotic branch
$\Veff\to 0$ providing the baseline for comparison.

\paragraph{$\Sthree\times\Sone$.}
As a conformally flat manifold ($C^2_\LC(\eta=V=0)=0$), the background Weyl
tensor's EC--Weyl coupling is inactive on the $\eta=V=0$ section.  However,
$\Sthree$ has a non-trivial geometry ($C^i_{jk}=4\varepsilon_{ijk}/r$), and
the spin-2 five-fold degeneracy
$\mu^2 = 128\pi^2 Lr/(3\kappa^2) - 16384\pi^2 L\alpha/(3r)$ and the spin-1
three-fold degeneracy $m^2_\omega = 16\pi^2 L^3/(\kappa^2 r)$ are established
as Theorem~4 (Palatini universality).  It also provides the foundation for
Pontryagin generation via the Pontryagin structure ($P\neq 0$ in MX, inherited
from paper~III).

\paragraph{$\Nilthree\times\Sone$.}
As the only non-conformally flat manifold among the three
($C^2_\LC=4/(3R^4)\neq 0$), this is the unique topology where the EC--Weyl
coupling generates a false vacuum for $\alpha<0$.  This false vacuum is
invisible in LC ($\alpha=0$) and first appears in the phase structure through
the EC extension.  The EC-specific new physics ($\gamma=1/2$ scaling, quintet
splitting, uniaxial mass splitting) established in Theorems~6--8 is understood
as the result of $\Nilthree$'s ``intermediate non-commutativity'' (between
$\Tthree$'s all-zero and $\Sthree$'s three-directional non-zero) being
physically realised as direction-dependent mode dictionaries.

This three-tier hierarchy is not a monotone increase in ``geometric complexity.''
Rather, $\Tthree$ (trivial) $\to$ $\Sthree$ (LC active, EC conformally
protected) $\to$ $\Nilthree$ (stage for EC-specific physics) represents a
division of roles that provides a comparative framework for homogeneous geometry.

\subsection{EC-Specific New Physics}
\label{sec:discussion:ecphysics}

We summarise the physical effects established in this paper that are specific to
EC (absent in LC).

\paragraph{(i) $\Nilthree$ false vacuum (Theorem~6).}
The $\Nilthree$ stable vacuum, which exhibits only a monotone potential in the
LC world, acquires a ``new stable phase'' through EC.  The scaling
$r_0\propto\sqrt{|\alpha|}$ exhibits a mechanism, non-trivial from the general
relativistic viewpoint, in which the strength of the Weyl coupling determines
the geometric size.

\paragraph{(ii) Quintet splitting (Theorem~7).}
The one-dimensional Lie algebra structure of $\Nilthree$ ($C^2_{01}$ only)
breaks the five-fold degeneracy of the spin-2 quintet around the EC false
vacuum to $5\to 0+2+1+1$.  In particular, the existence of zero mode $q_3$
shows the ``residual'' geometric symmetry of $\Nilthree$, with physically
different content from the LC isotropic point.

\paragraph{(iii) Uniaxial mass splitting (Theorem~8).}
The uniaxial selectivity in which only the non-commutative direction $\omega_2$
among the spin-1 triplet acquires a mass is a consequence of the asymmetric
internal structure of $\Nilthree$ being embedded in the kinematic matrix.  This
is the clearest manifestation of the direction-dependent mode dictionary in
$\Nilthree$.

\paragraph{(iv) $\alpha$-cubic coupling.}
In $\Sthree$, the EC-specific cubic coefficient
$\partial^3 V/\partial\alpha\partial\eta\partial V = -128\pi^2 V\eta r/3$ appears
(\S\ref{sec:s3:ec}).  This is a new nonlinear coupling added to the LC
NY-cubic channel ($= 4\pi^2 r^2$, invariant under EC), providing a cubic channel
in the homogeneous EFT.  On the other hand, since the Pontryagin density $P$ is
$\alpha$-independent ($\partial P/\partial\alpha=0$), this $\alpha$-cubic is
understood as a coupling located outside the Pontryagin structure.

\paragraph{(v) $\delta$-sector null result (Theorem~3).}
Complementarily, the $C_\delta$ defined for the formal auxiliary MIXING sector
is zero for all of $\Sthree, \Tthree, \Nilthree$, and $\partial_\alpha C_\delta=0$.
Therefore the new homogeneous cubic structure activated by the EC extension is
confined to the physical $\eta$-$V$ channel, and the auxiliary $\delta$-sector
remains null.

\subsection{What Is Not Claimed}
\label{sec:discussion:notclaimed}

The following are outside the scope of this paper.

\begin{itemize}
  \item Full stability including inhomogeneous modes beyond the minisuperspace.
  \item The claim that the full 4D deformed geometry forms a new Lie group
        manifold.
  \item Complete action generation including arbitrary non-invariant polynomials.
  \item A fully general homogeneous torsion ansatz including irreducible tensor
        torsion.  (This paper is restricted to the axial + vector-trace torsion
        ansatz adopted in papers~I--III.)
\end{itemize}

These are not in contradiction with the results of this paper; they are
consciously delineated as outside the stated scope.

\subsection{Toward Future Extensions}
\label{sec:discussion:future}

The homogeneous three-topology bulk sector established in this paper provides a
closed comparison framework in itself, while making explicit how different mode
dictionaries and protection structures appear in each topology.  The following
directions extend this framework naturally, while preserving the comparison
structure, toward local degrees of freedom, boundary responses, and physical
interpretation.

\begin{enumerate}
  \item Introduction of inhomogeneous perturbations containing local degrees of
        freedom.
  \item Problem settings including boundary surfaces and interfaces.
  \item Connection between homogeneous bulk modes and localised modes.
  \item Bridge to physical observables including thermal and Lorentzian
        interpretations.
\end{enumerate}
